% ============================================================
%  Paste this file into your Overleaf project.
%  Required packages (add to your main .tex preamble):
%    \usepackage{booktabs}
%    \usepackage{multirow}
%    \usepackage{amsmath}
%    \usepackage{algorithm}
%    \usepackage{algpseudocode}
% ============================================================

% ============================================================
\section{Introduction}
% ============================================================

Hong Kong's insurance market is one of the densest in Asia, with over 160
licensed insurers publishing product disclosure documents that mix Chinese and
English terminology, actuarial tables, and legally operative clause text.
Verifying a specific benefit, exclusion, or renewal condition requires locating
the right clause across dozens of PDFs---slow and inaccessible to most
policyholders.
Keyword search breaks down on paraphrased queries; general-purpose LLMs fail
differently, hallucinating policy details, ignoring questions the evidence
cannot support, and producing free-form prose with no audit trail.

The IRAG project tackles this through two successive development phases.
The first phase, carried out in a prior semester, established an initial
Milvus-based RAG prototype: a retrieval pipeline over Chinese insurance text
and a simple table-embedding scheme using a Google text embedding model
designed for English tables.
This prototype identified three concrete limitations.
The knowledge base was narrow in scale and insurer coverage.
The English-oriented table embedder performed poorly on Chinese insurance
tables and represented each table as a single coarse vector, making it
difficult to pinpoint the specific row or clause needed to answer a question.
Pilot fine-tuning experiments---a Chain-of-Thought run and a conventional
insurance-style answer run---showed promise but were not yet integrated with
the retrieval system.

Building on this foundation, the current semester addresses both limitations
in parallel.
On the retrieval side, the team expanded the knowledge base to 247 Hong Kong
insurance product PDFs across nine insurers, replaced the table embedder with
a fine-grained three-modal chunking scheme (full table, columns, rows), added
cross-encoder reranking, and deployed a FastAPI service with a Vue~3 chat
interface.
On the generation side, this paper describes the fine-tuning pipeline designed
to make the model safe to deploy in this RAG context.
We train Qwen3-4B~\cite{qwen3} through three successive LoRA stages: clause
reading comprehension, structured database reasoning via SQL chain-of-thought,
and Retrieval-Augmented Fine-Tuning (RAFT)~\cite{zhang2024raft} that teaches
the model to reason over retrieved chunks and refuse when those chunks do not
contain the answer.
We evaluate the result at three levels---offline behavioral tests against an
untuned baseline, a grounded attribution judge designed to catch hallucinations
that standard LLM-as-judge misses, and end-to-end smoke tests on the live
pipeline---to get a clear picture of what the fine-tuning fixes and where gaps
remain.

The rest of the paper is structured as follows.
Section~\ref{sec:related} reviews related work.
Section~\ref{sec:method} describes the three-stage training pipeline and
the RAG infrastructure.
Section~\ref{sec:exp} presents evaluation results.
Section~\ref{sec:discussion} discusses limitations and future directions.

% ============================================================
\section{Methodology}
\label{sec:method}
% ============================================================

\subsection{Base Model and Training Setup}

We use Qwen3-4B~\cite{qwen3} as the base language model---a 4-billion-parameter
transformer pre-trained on multilingual corpora with native Chinese support.
Hong Kong insurance clause text freely mixes Chinese and English terminology,
so bilingual coverage is a requirement.
All fine-tuning is managed by Axolotl with LoRA~\cite{hu2022lora} throughout
all three stages.
Adapters are attached to six projection layers per transformer block
($\mathbf{q}$, $\mathbf{k}$, $\mathbf{v}$, $\mathbf{o}$, \texttt{up\_proj},
\texttt{down\_proj}), with rank $r{=}32$, scaling factor $\alpha{=}64$, and
dropout 0.05.
Training runs on a single 16\,GB VRAM GPU (Blackwell architecture) in
\texttt{bfloat16} with gradient checkpointing enabled.
Flash Attention and Flex Attention are incompatible with this GPU generation
and are disabled.

\subsection{Three-Stage LoRA Fine-Tuning}

The pipeline uses three successive fine-tuning stages, each with a distinct
training objective and each initialised from the previous stage's LoRA
checkpoint.
Table~\ref{tab:stages} summarises the key hyperparameters.

\begin{table}[htbp]
\centering
\caption{Training configuration for the three fine-tuning stages.}
\label{tab:stages}
\small
\begin{tabular}{@{}lcccc@{}}
\toprule
Stage & seq\,len & lr & Epochs & Batch \\
\midrule
1 -- Clause   & 1024$^\dagger$ & $5{\times}10^{-5}$ & 1 & 2 \\
2 -- DB-CoT   & 1024           & $5{\times}10^{-5}$ & 4 & 4 \\
3 -- RAFT     & 1536           & $2{\times}10^{-5}$ & 3 & 4 \\
\bottomrule
\end{tabular}\\[2pt]
{\footnotesize
$\dagger$ Sliding window, 256-token overlap (Stage~1 only).
Batch = micro-batch $\times$ gradient accumulation steps.}
\end{table}

\noindent\textbf{Stage~1: Clause reading comprehension.}
Stage~1 trains in completion format on two insurance clause QA datasets---an
objective set and a subjective set at 40\% each---supplemented by a structured
database QA set at 20\%.
A 256-token sliding window handles contracts that exceed the 1024-token context
limit.
One epoch anchors the model's answers in the provided clause text and
establishes the expected output format.

\noindent\textbf{Stage~2: Structured database reasoning via SQL-CoT.}
The primary objective of Stage~2 is to teach the model \emph{structured,
evidence-grounded reasoning}: given a natural-language question and a
database schema, it should select the relevant fields, construct a step-by-step
SQL chain-of-thought (\texttt{<FieldSelection>} $\to$ \texttt{<SQLCoT>}),
and ground its final answer strictly in the query result rather than
parametric knowledge.
Training uses chat-format SQL-CoT samples derived from structured insurance
tables, with sample packing enabled.
Stage~2 training data contains only \emph{positive} reasoning
examples---question, schema, and correct answer---with no abstention examples.
The model therefore acquires structured reasoning but is not trained to refuse;
that capability is left entirely to Stage~3.
Four epochs give the model sufficient exposure to tabular reasoning without
displacing the clause-reading skills from Stage~1.

\noindent\textbf{Stage~3: RAFT fine-tuning.}
Stage~3 introduces retrieval-awareness and abstention on top of the
Stage~2 checkpoint.
Training uses the RAFT dataset described in Section~\ref{sec:raft_data}.
The learning rate drops to $2{\times}10^{-5}$ to limit forgetting of earlier
capabilities.
Sequence length grows to 1536 tokens, covering the 95th-percentile sample
length of 1472 tokens; only 3.8\% of samples require truncation.

\subsection{RAFT Training Data Construction}
\label{sec:raft_data}

We build the RAFT dataset from existing clause QA annotations and Stage~2 SQL
samples using a custom generator script (\texttt{generate\_raft.py}).
The generator produces approximately 1,995 training samples (after the
keyword-overlap filter removes RAG-positive candidates where key question
entities cannot be found in the retrieved chunk), distributed across three
categories at approximate target proportions of 52.5\,:\,31.5\,:\,15
(values are design targets; the accepted sample count varies with filtering
and is not guaranteed to sum to exactly 100\,\%):

\begin{enumerate}
  \item \textbf{RAG-positive.}
    Each sample pairs a question with 1--3 retrieved chunks, one of which
    contains the correct evidence.
    A keyword-overlap filter checks that key entities from the question
    appear in that chunk; candidates that fail this check are dropped.
    The model is trained to output a \texttt{<Thought>} reasoning block
    followed by a structured answer in three fields: \emph{answer},
    \emph{evidence}, and \emph{explanation}.

  \item \textbf{RAG-distractor.}
    The question is paired with two chunks drawn from unrelated contracts,
    so no chunk contains the correct answer.
    The model is trained to detect the mismatch and produce a standardised
    refusal rather than fabricating a response.

  \item \textbf{DB-CoT passthrough.}
    SQL-CoT samples from Stage~2 are included verbatim to prevent the model
    from losing structured-data reasoning during Stage~3 training.
\end{enumerate}

Questions fall into six types (definitional, conditional, judgement,
calculation, comparative, procedural), detected by keyword matching;
each type maps to a reasoning template that is injected into the
\texttt{<Thought>} block at generation time.

\noindent\textbf{Filtering parametric knowledge leakage.}
Stage~2 SQL-CoT training data occasionally elicits outputs that include
``专有名词解释'' (terminology explanation) paragraphs---free-form definitions
drawn from the model's pre-trained knowledge rather than retrieved evidence.
These paragraphs are harmless in a structured-database context where the task
is SQL generation, but they would constitute a grounding violation in RAFT
training: the model would learn to answer from memory rather than from
the provided chunks.
A post-processing step (\texttt{strip\_zhuanye()}) therefore removes any such
paragraph from RAFT training targets before they are written to disk.
The pattern (\texttt{专有名词*解释}) does not match clause text or SQL syntax,
so no evidence-bearing content is removed.

The system prompt instructs the model to answer strictly from retrieved
evidence and to refuse when that evidence is absent.

\subsection{Retrieval-Augmented Generation Pipeline}

The retrieval backend uses Milvus with \texttt{BAAI/bge-large-zh-v1.5}
embeddings (1024 dimensions, cosine similarity, IVF\_FLAT index with
$n_\text{list}{=}1024$).
The knowledge base indexes 37,571 chunks from 247 Hong Kong insurance product
PDFs spanning nine insurers: Prudential, Sun Life, AIA, Manulife, AXA, FWD,
BOC, Hang Seng, and HSBC.
Each PDF table is decomposed into three modality chunks---full table, columns,
and rows---to improve retrieval of tabular facts.
At query time, a FastAPI server (port 8000) retrieves the top-8 candidates
from Milvus, re-ranks them with \texttt{BAAI/bge-reranker-base} to sharpen
relevance ordering, and forwards the result to the RAFT inference server
(port 8001), which produces a response grounded in the retrieved content.


% ============================================================
\section{Experiments}
\label{sec:exp}
% ============================================================

\subsection{Evaluation Protocol}

Evaluation covers three aspects of the IRAG system.

\noindent\textbf{Model capability.}
We define \emph{legalized} outputs as responses that (i) are grounded in
retrieved evidence, (ii) adopt a structured, auditable format, and
(iii) refuse when evidence is insufficient---criteria intended to align model
behavior with the accountability requirements of insurance advising.
We test whether fine-tuning reliably produces such outputs.
We compare Baseline (Qwen3-4B without fine-tuning) against Stage-3 (RAFT) on
a suite of offline tests covering three behavioral dimensions:
evidence grounding, abstention reliability, and structured output compliance.
We do not separately evaluate Stage-2, whose contribution to the final model
is validated indirectly through the grounding and compliance metrics of Stage-3.

\noindent\textbf{Knowledge base retrieval.}
Retrieval quality requires domain expertise to judge accurately; we therefore
do not report a standalone retrieval benchmark.
The end-to-end evaluation covers retrieval as one of four scoring dimensions.

\noindent\textbf{End-to-end pipeline.}
We run the complete IRAG pipeline---Milvus retrieval plus Stage-3
generation---on real Hong Kong insurance questions and score system-level outputs.

All offline judgements use DeepSeek-Chat~\cite{deepseek} at temperature 0.1;
end-to-end judgements use temperature 0.0.
Model generation runs at temperature 0.2, top-$p$ 0.9, maximum 400 new tokens.
Models are loaded sequentially to stay within 16\,GB VRAM.
Because LLM judges routinely miss implicit hallucination, we also run a grounded
judge that traces each factual claim against the source text before scoring
(Section~\ref{sec:abstention}).

\subsection{Model Capability: Legalized Output}

\subsubsection{Evidence Grounding}

\noindent\textbf{Answer compliance.}
We score 30 clause QA output pairs on three compliance dimensions: answer
compliance (precise use of insurance terminology), clause fidelity (direct
citation of clause text rather than paraphrase), and response formality
(3\,pts each; 15\,pts total).
Both models receive the deployment system prompt, which includes format
directives (\texttt{<Thought>} block and structured answer fields).
Baseline ignores these directives and produces free-form prose; Stage-3
consistently applies the structured format as trained.
The response formality dimension therefore captures a real behavioral
difference, not a prompt-induced artifact.

\begin{table}[htbp]
\centering
\caption{Answer compliance results (avg\,/\,3 per dimension).
Bold indicates the better score.}
\label{tab:compliance}
\small
\begin{tabular}{@{}lcc@{}}
\toprule
Dimension & Baseline & Stage-3 \\
\midrule
Answer compliance   & 2.93 & \textbf{3.47} \\
Clause fidelity     & 2.57 & \textbf{3.43} \\
Response formality  & \textbf{3.80} & 3.70 \\
\midrule
Total               & 62.0\% & \textbf{70.7\%} \\
\bottomrule
\end{tabular}\\[4pt]
{\footnotesize \textit{Rubric} (per dimension, 3\,pts max): 3\,=\,fully meets
criterion; 2\,=\,partially meets (e.g.\ paraphrase instead of direct citation);
1\,=\,does not meet. \emph{Answer compliance}: precision of domain terminology.
\emph{Clause fidelity}: extent of verbatim clause citation.
\emph{Response formality}: professional register throughout.}
\end{table}

Stage-3 leads overall (70.7\% vs.\ 62.0\%), with the largest margin on clause
fidelity (+0.86): Stage-3 tends to quote clause text directly, while Baseline
paraphrases or summarises.
The gain likely comes from Stage~3 RAFT training, which rewards quoting clause
text directly; SQL-CoT training in Stage~2 may also play a role by conditioning
the model to reference specific fields rather than paraphrase.
Which contribution dominates is outside the scope of this work.

\noindent\textbf{Fair clause evaluation.}
We sample the same 30 questions from \texttt{eval/data/clause.jsonl} for a
direct QA accuracy comparison.
Baseline receives a single-chunk prompt with a clause-focused system prompt;
Stage-3 receives a two-chunk RAFT-format input (one correct evidence chunk
plus one distractor) with the deployment system prompt.
The judge sees each model's actual input, so Stage-3 is not penalized for
responding to a richer context than Baseline.
Five dimensions are scored (correctness, evidence use, no-hallucination,
structure, fidelity; 5\,pts each; 25\,pts total).

\begin{table}[htbp]
\centering
\caption{Fair clause evaluation results (avg\,/\,5 per dimension).
Bold indicates the better score.}
\label{tab:clause}
\small
\begin{tabular}{@{}lcc@{}}
\toprule
Dimension & Baseline & Stage-3 \\
\midrule
Correctness       & \textbf{3.43} & 2.53 \\
Evidence use      & 3.53          & \textbf{4.07} \\
No-hallucination  & \textbf{4.90} & 4.17 \\
Structure         & \textbf{4.70} & 4.10 \\
Fidelity          & \textbf{4.13} & 3.57 \\
\midrule
Total             & \textbf{82.8\%} & 73.7\% \\
\bottomrule
\end{tabular}
\end{table}

Stage-3 leads on evidence use (4.07 vs.\ 3.53), the dimension most directly
measuring whether the answer is derived from retrieved text.
The overall gap (82.8\% vs.\ 73.7\%) is concentrated in correctness and
structure.
We attribute this primarily to a format mismatch: Stage-3 consistently applies
its RAFT output template---\texttt{<Thought>} block followed by structured
answer fields---even where a direct answer would score higher under the clause
rubric.
Some loss of direct-answer fluency cannot be ruled out; an ablation without
the format constraint would be needed to separate the two effects.

\subsubsection{Abstention Reliability}
\label{sec:abstention}

\noindent\textbf{Hallucination suppression.}
We test 10 questions where the provided evidence is deliberately insufficient
to answer completely.
Stage-3 is expected to refuse and provide an advisory response; Baseline has
no such training.
Both models receive the deployment system prompt; the evidence provided is
deliberately insufficient to fully answer each question.

\begin{table}[htbp]
\centering
\caption{Hallucination suppression results (avg\,/\,5 per dimension).
Bold indicates the better score.}
\label{tab:hall}
\small
\begin{tabular}{@{}lcc@{}}
\toprule
Dimension & Baseline & Stage-3 \\
\midrule
Correctness       & \textbf{4.00} & 3.50 \\
Evidence use      & \textbf{4.00} & 3.30 \\
No-hallucination  & \textbf{4.50} & 3.60 \\
Structure         & \textbf{4.60} & 4.10 \\
Fidelity          & \textbf{4.20} & 3.50 \\
\midrule
Total             & \textbf{85.2\%} & 72.0\% \\
\bottomrule
\end{tabular}
\end{table}

Baseline leads across all five dimensions under the standard judge
(85.2\% vs.\ 72.0\%), but the gap has two separate causes.
When Stage-3 refuses, it produces no direct answer; the judge penalises this
on correctness and fidelity regardless of whether the refusal is warranted.
Additionally, Stage-3 refusals often include short advisory statements
explaining what kind of evidence would be needed; the judge treats these as
factual claims, pulling the no-hallucination score down (3.60 vs.\ 4.50).
On the other side, Baseline's high no-hallucination score (4.50) is partly
inflated: Baseline answers confidently even when its response goes beyond the
provided evidence, but the standard judge does not consistently catch this.
The next test addresses this asymmetry directly.

\noindent\textbf{Grounded judge correction.}
We re-score the same outputs with a two-stage judge.
First, each factual claim in a response is checked against the source text;
any claim without direct textual support forces \texttt{no\_hallucination=0}
and \texttt{evidence\_use=0}.
Standard five-dimension scoring then proceeds on top of the attribution result.

\begin{table}[htbp]
\centering
\caption{Standard vs.\ grounded judge on hallucination suppression.
\textcolor{red}{Red}: score decreases; \textcolor{blue}{blue}: score increases.}
\label{tab:grounded}
\small
\begin{tabular}{@{}lccc@{}}
\toprule
Model & Standard & Grounded & $\Delta$ \\
\midrule
Baseline & 85.2\% & 80.4\% & \textcolor{red}{$-$4.8\%} \\
Stage-3  & 72.0\% & 74.0\% & \textcolor{blue}{$+$2.0\%} \\
\midrule
Gap & 13.2\% & \textbf{6.4\%} & \\
\bottomrule
\end{tabular}
\end{table}

Attribution checking narrows the gap from 13.2\% to 6.4\%.
Baseline loses points on 3 of 10 questions where it gives authoritative-sounding
answers that cannot be traced to the provided text.
Stage-3 gains slightly: once the attribution check is in place, refusals
correctly receive zero hallucination credit.
The remaining 6.4\% gap likely reflects implicit hallucination that neither
judge fully surfaces, a limitation of automated evaluation for this class of
behavior.

\noindent\textbf{Distractor robustness.}
The harder abstention test gives the model two chunks from \emph{unrelated}
contracts; the correct response is to recognise the mismatch and refuse.
We evaluate 15 questions on three dimensions (rejection correctness,
no-hallucination, explanation quality; 5\,pts each; 15\,pts total),
using the deployment system prompt for both models.
Stage-3 scores 31.1\% and Baseline 43.1\%.
Stage-3 trails Baseline here: RAG-distractor samples account for only 31.5\%
of the RAFT training set, and the dataset lacks a ``partially relevant''
category---cases where retrieved content is topically close but does not
contain the specific clause---that would let the model learn a finer refusal
threshold.
We discuss this further in Section~\ref{sec:discussion}.

\subsubsection{Structured Output Compliance}

We parse all 15 Stage-3 responses from the end-to-end smoke test using
regular expressions to measure format compliance under live traffic.

\begin{table}[htbp]
\centering
\caption{Structured output compliance across 15 live pipeline outputs.
Baseline produces no structured elements by construction (no format training).}
\label{tab:format}
\small
\begin{tabular}{@{}lcc@{}}
\toprule
Format element & Baseline & Stage-3 \\
\midrule
\texttt{<Thought>} block present      & 0\,/\,15 (0\%) & 15\,/\,15 (100\%) \\
\texttt{[答案]} field present          & 0\,/\,15 (0\%) & 15\,/\,15 (100\%) \\
Full 3-field structure (refusal only) & 0\,/\,2\hphantom{0} (0\%) & \phantom{0}2\,/\,2\hphantom{0} (100\%) \\
\bottomrule
\end{tabular}
\end{table}

Baseline produces none of the required structural elements, as it received no
format-specific training.
Stage-3 achieves 100\% compliance on all three elements across 15 outputs,
applying an adaptive format: answering questions emit a \texttt{<Thought>}
block followed by \texttt{[答案]} only, while refusals switch to the full
three-field structure (\texttt{[答案]/[证据]/[解释说明]}).
Stage-3 applies the format consistently across live traffic.

\subsection{End-to-End Pipeline Evaluation}
\label{sec:e2e}

We submit real Hong Kong insurance questions to the live IRAG
pipeline---Milvus retrieval (top-8 chunks) followed by Stage-3
generation---and score each response on four dimensions: retrieval relevance,
faithfulness, rejection quality, and completeness (3\,pts each; 12\,pts total),
judged by DeepSeek-Chat.
A response scoring 12 is marked \textsc{pass}; anything lower is
\textsc{fail}.

\begin{table}[htbp]
\centering
\caption{End-to-end smoke test results.}
\label{tab:smoke}
\small
\begin{tabular}{@{}lcccc@{}}
\toprule
Set & \#Q & PASS & FAIL & Avg\,/\,12 \\
\midrule
v4 & 10 & 7 (70\%) & 3 (30\%) & 7.90 \\
v5 & 15 & 9 (60\%) & 6 (40\%) & 8.07 \\
\bottomrule
\end{tabular}
\end{table}

The v5 set extends coverage from 6 to 9 insurers.
The 60\% pass rate is lower than v4's 70\%, but the per-question average is
slightly higher (8.07 vs.\ 7.90): questions with imperfect retrieval still earn
partial credit on faithfulness and rejection.

\noindent\textbf{Failure taxonomy.}
The six v5 failures fall into three layers (Table~\ref{tab:failures}).

\begin{table}[htbp]
\centering
\caption{Failure taxonomy for the v5 smoke test (6 failures).}
\label{tab:failures}
\small
\begin{tabular}{@{}p{1.3cm}p{1.2cm}p{4.5cm}@{}}
\toprule
Layer & Questions & Root cause \\
\midrule
Retrieval & Q3, Q7, Q9 & Chunks are on-topic but lack the specific numerical values required \\
KB coverage & Q12 & Hang Seng life product absent from the index \\
Generation & Q11, Q13 & Evidence spans multiple chunks near the context limit; answer is incomplete \\
\bottomrule
\end{tabular}
\end{table}

Among the passing questions, Q15 (the refusal question) scores 6\,/\,12---it
declines correctly without drawing on parametric knowledge, showing that
Stage~3's abstention training holds under live conditions.
Note that the 12\,/\,12 pass threshold is calibrated for answerable questions;
Q15's score of 6\,/\,12 represents a correct refusal rather than a failure in
the intended sense---it is counted as \textsc{fail} in the aggregate statistics
to avoid inflating the pass rate, but interpreted as a positive outcome;
excluding Q15, the pass rate on answerable questions is 9\,/\,14 (64.3\%).
Retrieval-layer failures are addressable at the system level---higher top-$k$
or a stronger re-ranker---without retraining the model.
The two generation-layer failures reflect the 1536-token context limit: when
relevant information is spread across several chunks, the answer is truncated.
The KB coverage gap for Q12 requires adding the missing insurer documents to
the Milvus index.


% ============================================================
\section{Discussion}
\label{sec:discussion}
% ============================================================

\subsection{Retrieval as the Performance Ceiling}

The end-to-end smoke test (Section~\ref{sec:e2e}) shows that three of the six
v5 failures originate in the retrieval layer: the correct document is indexed,
but the retrieved chunks describe the right topic without including the specific
numerical value or clause condition the question requires.
This is not a model failure---Stage-3 correctly reports uncertainty when the
chunks are uninformative---but it means that retrieval quality sets a hard
ceiling on what fine-tuning can achieve.
The current setup retrieves top-8 candidates and re-ranks with a cross-encoder;
raising top-$k$ or adding a query expansion step are the most direct routes to
closing this gap without retraining the model.

The knowledge base also has uneven insurer coverage.
Q12 failed entirely because Hang Seng life products were absent from the index.
With 37,571 chunks from 247 PDFs, coverage is broad but not complete;
any production deployment would need a process for identifying and filling
coverage gaps as new products are issued.

\subsection{Fine-Tuning Trade-offs: Grounding vs.\ Direct Accuracy}

The fair clause evaluation reveals a tension between evidence grounding and
raw QA accuracy.
Stage-3 leads on evidence use (4.07 vs.\ 3.53) but trails on total score
(73.7\% vs.\ 82.8\%), with the gap concentrated in correctness and structure.
The most likely cause is a format mismatch: Stage-3 consistently applies its
RAFT output template even in settings where a direct answer would score higher
under the clause rubric.
An ablation that removes the format constraint at inference time would separate
this effect from any genuine loss of answer fluency, but we leave that to
future work.

The compliance evaluation tells a complementary story.
Stage-3 leads overall (70.7\% vs.\ 62.0\%), with the largest individual gain
on clause fidelity (+0.86 on a 3-point scale).
Together, the two evaluations suggest that fine-tuning shifts the model's
behavior toward citation and away from paraphrase---a useful property in
an auditable insurance QA context, even if it costs some points on
correctness-focused rubrics.

\subsection{Abstention: Capability Transfer but Unresolved Trade-off}

Stage-3 transfers its abstention capability to the live pipeline---Q15's
correct refusal (6\,/\,12) in the smoke test confirms this---but the offline
distractor test exposes a persistent weakness: Stage-3 scores 31.1\% against
Baseline's 43.1\% on questions where both chunks come from unrelated contracts.

Two factors explain this.
RAG-distractor samples make up only 31.5\% of the RAFT training set, giving
the model limited exposure to the harder rejection cases.
More fundamentally, the training data has no ``partially relevant'' category:
cases where retrieved content is topically close but does not contain the
specific clause.
The model has learned a binary distinction (relevant vs.\ clearly irrelevant)
that does not generalise to the more ambiguous cases that appear in practice.

We also explored three system prompt configurations during development.
A strict refusal prompt (v1) maximised distractor rejection but degraded
the advisory quality of refusal responses.
A graded response prompt (v2) improved hallucination suppression but caused
distractor rejection to collapse.
The deployment prompt (v3) sits between the two, but the trade-off cannot be
resolved by prompt tuning alone.
The fundamental fix is adding a partially-relevant training category so
the model can learn a finer rejection threshold.

\subsection{Evaluation Methodology: LLM-as-Judge Limitations}

The grounded judge experiment (Section~\ref{sec:abstention}) makes explicit
what automated evaluation misses.
Under the standard judge, Baseline outscores Stage-3 by 13.2 percentage points
on the hallucination suppression task.
After attribution checking---tracing each factual claim to the source
text---the gap narrows to 6.4 points.
Baseline loses credit because three of its ten responses contain
authoritative-sounding claims that cannot be located in the provided evidence;
Stage-3 gains slightly because its refusals, once evaluated with the
attribution requirement applied, correctly receive zero hallucination
penalties.

The remaining 6.4-point gap likely reflects implicit hallucination that neither
judge can surface: confident paraphrases that are plausible but untraceable.
This is a general limitation of automated evaluation for this class of
behavior, and it motivates caution in interpreting all LLM-as-judge results
in this paper.
Human evaluation on a larger question set would give a more reliable picture,
particularly for the hallucination and distractor tasks where sample sizes
(10 and 15 questions respectively) are small enough that individual questions
have a noticeable effect on the aggregate.

\subsection{System-Level Considerations}

One evaluation in this paper was invalidated by a system-level issue rather
than a model failure.
The noise filter test---designed to measure whether Stage-3 selects the correct
chunk from a mix of relevant and irrelevant inputs---produced misleading results
because 10 of the 30 evaluation chunks had malformed metadata (contract name
listed as ``unknown'').
Stage-3 treated these as unverifiable sources and refused to answer, depressing
its correctness score well below Baseline.
In the production pipeline, Milvus extracts contract names from PDF metadata at
ingestion time, so this problem does not arise.
The episode illustrates a general point: a model trained to be conservative
about sourcing is more sensitive to data quality issues upstream, and
evaluation pipelines need to match the data conditions of the deployment
environment.

\subsection{Future Directions}

Several concrete improvements follow from the analysis above.

At the retrieval level, hybrid dense--sparse retrieval (combining vector search
with BM25) may improve recall for queries containing rare insurance terms that
the embedding model has not seen frequently.
Higher top-$k$ retrieval with more aggressive reranking would reduce the
chance that a needed clause is simply not in the candidate set.

At the training level, the most pressing need is a partially-relevant
distractor category in the RAFT dataset: samples where retrieved content
overlaps topically but does not contain the specific clause, trained to produce
a calibrated partial-refusal response rather than either answering or fully
refusing.
Increasing the distractor ratio from 31.5\% would also help, at the cost of
reducing positive-sample exposure.

At the evaluation level, the grounded judge described in this paper is a step
toward more reliable hallucination measurement, but it still relies on an LLM
to perform attribution checking.
A deterministic attribution method---exact string matching or a specialised
entailment model trained on insurance clause text---would remove this source of
noise.
Larger test sets for the hallucination and distractor tasks are also needed
before the results can be reported with confidence intervals narrow enough to
support strong conclusions.
